\documentclass{report}
\usepackage[utf8]{inputenc}
\usepackage[T1]{fontenc}
\usepackage{geometry}
\usepackage{graphicx}
\usepackage{float}
\usepackage{hyperref}
\usepackage[english,indonesian]{babel}

\title{Optimasi Pemrosesan Template LaTeX Menggunakan Pendekatan Berbasis AI}
\author{Academic Support Agent}
\date{\today}

\begin{document}
\maketitle

\begin{abstract}

\end{abstract}

\chapter{PENDAHULUAN}


\chapter{TINJAUAN PUSTAKA}


% Original content (commented out):
% % === No modificar éstos parámetros === 

\documentclass[letterpaper, 12pt]{report}
\usepackage[utf8]{inputenc}
\usepackage[english, spanish]{babel}
\usepackage{fullpage} % changes the margin
\usepackage{graphicx} 
\usepackage{enumitem} 
\usepackage{chngcntr}
\counterwithin{figure}{section}
\renewcommand{\thesection}{\arabic{section}} 
\renewcommand{\thesubsection}{\thesection.\arabic{subsection}}
\renewcommand{\baselinestretch}{1.5}


\begin{document}
% Hoja de portada, únicamente editar nombres y códigos

\begin{titlepage}
	\centering
	\includegraphics[width=0.3\textwidth]{logo_utb.png}\par\vspace{1cm}
	{\scshape\LARGE Universidad Tecnológica de Bolívar \par}
	\vspace{1cm}
	{\scshape\Large Programa de Ingeniería de Sistemas y Computación \par}
	\vspace{.2cm}
    {\scshape\Large Nombre del curso \par}
	\vspace{1cm}
	{\Large\bfseries Nombre de la actividad\par}
	\vspace{1cm}
	{\itshape Estudiante 1, T000XXXX \par}
    {\itshape Estudiante 3, T000XXXX \par}
    {\itshape Estudiante 3, T000XXXX \par}
    {\itshape Estudiante 4, T000XXXX \par}
	\vfill
	Prof. Nombre Apellido
	\vfill
	{\large \today\par}
\end{titlepage}

\newpage

\section{Primera parte}
En ésta sección se deben explicar en detalle los aspectos m\'as relevantes de la práctica realizada, aprendizaje obtenido y posibles usos y aplicaciones del montaje realizado. 

\newpage

\section{Segunda parte}
En ésta sección se deben explicar en detalle los aspectos m\'as relevantes de la práctica realizada, aprendizaje obtenido y posibles usos y aplicaciones del montaje realizado. 

\newpage

\section{Conclusiones}
En ésta sección se deben explicar en detalle los aspectos m\'as relevantes de la práctica realizada, aprendizaje obtenido y posibles usos y aplicaciones del montaje realizado. 

\begin{thebibliography}{9}
\bibitem{rashid} Rashid, M. H. (2001). Power Electronics Handbook. Power Electronics Handbook (Third Edition).http://dx.doi.org/10.1016/B978-0-12-382036-5.00036-7
\end{thebibliography}

\end{document}


\end{document}